\bookchapter{第零章：噫，始之前}

正式探索前，且先論基本規則，以保後述汝皆知。勿急，磨刀不誤砍柴工也。

\booksection{一、此書能助汝何爲？（與不能爲之事）}

其能助汝系統性通曉手機何以用，明功能背後之理，且無論汝乃安卓黨抑或蘋果派，多數思路皆通用。其將引汝動手練習，使汝自“會用”升至“懂用”。

其能助汝系統性通曉手機何以用，明功能背後之理，且無論汝乃安卓黨抑或蘋果派，多數思路皆通用。其將引汝動手練習，使汝自“會用”升至“懂用”。

\booksection{二、吾等之“暗號”釋}

爲說得更明，書中將用些許小標記，此等標記猶如吾等間之默契暗號：

\begin{itemize}
  \item \textbf{[按鈕名]}：表屏幕上具體之按鈕或App之名，譬如[設置]、[確定]。
  \item \textbf{>}：表操作之序，譬如"點 [設置] > [聲音]"，即先點設置，再點聲音。
  \item \textbf{加粗}：至要！或爲關鍵概念或安全警告，請務必留意。
  \item \emph{斜體}：些許補充之小想或比喻，助汝更善理解。
\end{itemize}

\booksection{三、安卓與iOS，皆論之}

汝猜何如，有些事物稱謂不同，然所爲之事一也。譬如：應用（安卓）與 App（iOS），設置（皆稱設置），文件管理（安卓）與 文件（iOS）。知其爲同一事即可，勿爲名詞所惑。

此外，今安卓陣營中尚有華爲之鴻蒙系統，其有自身特色，後吾等將專論之。總之，不論汝用何手機，此書皆能助汝。

\booksection{四、手機容量：何以買128G，可用卻不及110G？}

容量單位衆所知：1TB = 1024GB，1GB = 1024MB，1MB = 1024KB……然此處有一坑：硬盤廠商通常按1000進制算，而系統按1024進制算，故汝所購之“128G”硬盤，實際容量將“縮水”。

然汝手機容量縮水之原因主要有三：

\begin{enumerate}
  \item \textbf{系統自佔其地}：操作系統猶如房屋之承重牆，必存。今之系統動輒佔用10-20GB，此甚正常。
  \item \textbf{廠商預裝之應用}：新手機中自帶之App亦佔空間，有能卸者，有不能者。
  \item \textbf{分區與格式化損耗}：技術之故，猶如書包之襯布亦佔體積。
\end{enumerate}

\begin{tcbinfo}{買手機之小建議}
買手機之小建議

\end{tcbinfo}

\begin{tcbpractice}{來動手：觀汝手機之"身份證"}
來動手：觀汝手機之“身份證”

\end{tcbpractice}

\booksection{五、汝之手機乃一座“數字資產庫”}

始具體知識前，吾等先立一核心認知：\textbf{手機早已不唯通信工具，乃吾等最重要之數字資產庫}。聯繫人、聊天記錄、家庭照片、工作文檔、銀行App、付款碼、健康數據、電子證件……此等數據相加，價值往往逾手機本身。

\begin{tcbinfo}{此何以至要？}
此何以至要？\textbf{更恐爲多年積累之不可再生數據}——孩子之成長照片、過世親人之語音、未備份之重要合同。科學地整理與防護習慣，能自三維護汝：

\begin{itemize}
  \item \textbf{提升運行效率}：規範之文件分類使系統查找更速，手機長期不卡頓。
  \item \textbf{延使用壽命}：合理之存儲管理減芯片讀寫之負，電池與存儲皆更耐用。
  \item \textbf{保數據安全}：分層管理 + 定期備份 + 權限收斂，使"數據災難"遠離汝。
\end{itemize}

\end{tcbinfo}

\booksubsection{吾等之核心理念：“預防爲主，防治結合”}

猶養生不得待病方就醫，\textbf{手機管理亦須自“日常習慣”抓起}。本書反覆強調者非某一高級技巧，乃此等樸素之原則：

\begin{itemize}
  \item \textbf{預防爲主}：裝 App 唯走官方商店、照片發前先脫敏、密碼不重複——此等"小動作"成本極低，然能擋絕大多數風險。
  \item \textbf{定期體檢}：每月一整理下載文件夾、每季一權限審計、每半年一全量備份——固定節奏比臨時救火有效百倍。
  \item \textbf{留有退路}：啓雲同步、記好賬號密碼、書一份"數字遺產清單"——確保意外發生時汝與家人皆能從容應對。
\end{itemize}
